\begin{derivation}[结构函数矩与局域算符矩阵元]
深度非弹性极限由 $\mathcal Q^2\gg\Lambda_{\rm had}^2$ 控制；Fourier 积分的主贡献来自 $|z^2|=O(\hbar^2/\mathcal Q^2)$ 的短距离区域。对任意选定的有限 OPE 截断集合 $\mathscr S$，先把余项定义为
\begin{equation*}
 \mathcal R_{\mu\nu}^{(\mathscr S)}(z)
 \equiv
 \mathcal T\{j_\mu(z)j_\nu(0)\}
 -\sum_{(N,i)\in\mathscr S}
 C_{\mu\nu,N}^{(i)}(z,\mu_{\rm F})
 \mathcal O_N^{(i)}(0,\mu_{\rm F}).
\end{equation*}
于是有限阶关系是严格等式
\begin{equation*}
 \mathcal T\{j_\mu(z)j_\nu(0)\}
 =
 \sum_{(N,i)\in\mathscr S}
 C_{\mu\nu,N}^{(i)}(z,\mu_{\rm F})
 \mathcal O_N^{(i)}(0,\mu_{\rm F})
 +\mathcal R_{\mu\nu}^{(\mathscr S)}(z).
\end{equation*}
OPE 的动力学内容是：按扭曲度与算符维数扩大 $\mathscr S$ 时，$\mathcal R_{\mu\nu}^{(\mathscr S)}$ 在上述短距离区域按相应的 $\Lambda_{\rm had}^2/\mathcal Q^2$ 幂次受抑；因此截断误差与所保留的扭曲度层级一一对应。
领先的非极化夸克算符是对称无迹、扭曲度为 $2$ 的张量
\begin{equation*}
 \mathcal O_{f,N}^{\mu_1\cdots\mu_N}
 =\bar\psi_f\gamma^{\{\mu_1}
 (\ii\hbar D^{\mu_2})\cdots
 (\ii\hbar D^{\mu_N\}})\psi_f
 -\text{迹项}.
\end{equation*}
每个 $\ii\hbar D^\mu$ 都具有物理动量量纲。Lorentz 协变性把前向矩阵元压缩为一个约化标量系数
\begin{equation*}
 \langle P|\mathcal O_{f,N}^{\mu_1\cdots\mu_N}|P\rangle
 =2A_{f,N}(\mu_{\rm F})
 P^{\{\mu_1}\cdots P^{\mu_N\}}
 -\text{迹项},
\end{equation*}
按此定义 $A_{f,N}$ 无量纲。

PDF 与这些局域算符之间的联系来自类光算符对 $\tau$ 的规范协变 Taylor 展开。先把 Wilson 线与被平移的场搬到同一端点：
\begin{equation*}
 \mathcal W_n(0,\tau)\psi_f[z(\tau)]
 =\sum_{r=0}^{\infty}
 \frac1{r!}
 \left(\frac{\hbar\tau}{P\cdot n}\right)^r
 (n\cdot D)^r\psi_f(0).
\end{equation*}
对厄米共轭双局域算符同样展开，便逐阶产生带协变导数的局域算符。另一方面
\begin{equation*}
 x^{N-1}\ee^{-\ii x\tau}
 =(\ii\partial_\tau)^{N-1}\ee^{-\ii x\tau}.
\end{equation*}
把 $\partial_\tau$ 通过分部积分移到 $\mathcal B_{f/H}$ 上，每次导数都带出 $\ii\hbar n\cdot D/(P\cdot n)$。因此 Mellin 矩与局域矩阵元一一对应：
\begin{equation*}
 \int_0^1\dd x\,x^{N-1}
 \left[f_{f/H}(x,\mu_{\rm F})+(-1)^Nf_{\bar f/H}(x,\mu_{\rm F})\right]
 =A_{f,N}(\mu_{\rm F}).
\end{equation*}
所以局域 OPE 的算符混合与 DGLAP 的 PDF 演化并不是两套独立理论；它们分别是同一类光算符重整化在 Mellin 空间与 $x$ 空间中的表示。
\end{derivation}

\begin{derivation}[从矩展开到共线卷积]
定义 Mellin 卷积

\begin{equation*}
 [A\otimes B](x)
 \equiv\int_x^1\frac{\dd z}{z}
 A(z)B\!\left(\frac{x}{z}\right).
\end{equation*}
它满足结合律：
\begin{align*}
 [A\otimes(B\otimes C)](x)
 &=\int_x^1\frac{\dd z}{z}A(z)
 \int_{x/z}^1\frac{\dd u}{u}B(u)
 C\!\left(\frac{x}{zu}\right),\\
 &=[(A\otimes B)\otimes C](x),
\end{align*}
第二行只交换积分次序并令乘积变量 $v=zu$。这个代数性质保证因子化反项可以在系数函数与 PDF 之间一致移动。

在 $d=4-2\epsilon$ 维正则化中，裸类光算符发生混合，写成卷积矩阵
\begin{equation*}
 f_i^{\rm bare}
 =\sum_j Z_{ij}^{\rm F}\otimes f_j(\mu_{\rm F}).
\end{equation*}
在 $\overline{\rm MS}$ 方案的领先阶，
\begin{equation*}
 Z_{ij}^{\rm F}(z)
 =\delta_{ij}\delta(1-z)
 +\frac{\alpha_s}{2\pi}
 \frac1{\bar\epsilon}
 P_{ij}^{(0)}(z)
 +\order(\alpha_s^2),
\end{equation*}
其中 $1/\bar\epsilon=1/\epsilon-\gamma_E+\ln4\pi$。部分子级系数函数中的相反共线极点由 $(Z^{\rm F})^{-1}$ 消去，留下有限的短距离系数 $C_{a,i}$。于是领先幂结构函数成为
\begin{equation*}
 \mathfrak F_a(x_B,\mathcal Q^2)
 =\sum_i
 \left[
 C_{a,i}\!\left(
 \frac{\mathcal Q^2}{\mu_{\rm F}^2},
 \frac{\mu_{\rm DR}^2}{\mu_{\rm F}^2},
 \alpha_s(\mu_{\rm DR})
 \right)
 \otimes f_{i/H}(\mu_{\rm F})
 \right](x_B)
 +\order\!\left(\frac{\Lambda_{\rm QCD}^2}{\mathcal Q^2}\right).
\end{equation*}
这就是部分子模型在 QCD 中的受控推广：$C_{a,i}$ 可由微扰理论计算，$f_{i/H}$ 是规范不变的非微扰强子矩阵元，而两者对人为因子化尺度的依赖必须互相抵消。
\end{derivation}

\begin{derivation}[夸克与胶子的相互作用顶角]
把协变导数
\[
D_\mu=\partial_\mu+\ii g_s\mathscr C_\mu^aT^a
\]
代入夸克动能项，先把自由项与相互作用项严格分开：
\begin{align*}
\bar\psi_f\ii\hbar c\gamma^\mu D_\mu\psi_f
&=\bar\psi_f\ii\hbar c\gamma^\mu\partial_\mu\psi_f
+\bar\psi_f\ii\hbar c\gamma^\mu
(\ii g_s\mathscr C_\mu^aT^a)\psi_f\\
&=\bar\psi_f\ii\hbar c\gamma^\mu\partial_\mu\psi_f
-\hbar c g_s\bar\psi_f\gamma^\mu T^a\psi_f\,\mathscr C_\mu^a.
\end{align*}
因此三场相互作用拉氏量就是
\[
\Lag_{q\bar q g}
=-\hbar c g_s\bar\psi_{f,i}\gamma^\mu(T^a)_{ij}\psi_{f,j}\mathscr C_\mu^a.
\]
Dyson 展开的一阶项为
\begin{align*}
S_{q\bar q g}^{(1)}
&=\frac{\ii}{\hbar c}\int\dd^4x\,\Lag_{q\bar q g}(x)\\
&=-\ii g_s\int\dd^4x\,
\bar\psi_{f,i}(x)\gamma^\mu(T^a)_{ij}\psi_{f,j}(x)\mathscr C_\mu^a(x).
\end{align*}
把三条场都作 Fourier 展开，局域乘积中的时空积分产生整体四动量守恒因子
\[
\int\dd^4x\,
\exp\!\left[-\frac{\ii}{\hbar}
(p_1+p_2+k)\cdot x\right]
=(2\pi\hbar)^4\delta^{(4)}(p_1+p_2+k),
\]
这里按图规则约定三条四动量全部流入顶角。剥去所有图共有的动量守恒 $\delta$ 函数以后，剩下的局域矩阵系数恰为
\[
{-\ii g_s\gamma^\mu(T^a)_{ij}}.
\]
因此一个 QCD 顶角的三部分来源可以逐项追溯：$-\ii g_s$ 来自 Dyson 指数与协变导数中的耦合，$\gamma^\mu$ 来自 Dirac 电流，$T^a$ 来自颜色联络在基本表示上的生成元。图规则只是把这三个作用量系数压缩成一个可重复使用的记号。
\end{derivation}

\begin{derivation}[不同味夸克散射的自旋与颜色平均]
从\eqnrefs{eq:qqprime-amplitude-dp10} 出发，初态自旋平均和颜色平均可分开写成
\begin{equation*}
 \overline{|\mathcal M|^2}
 =\frac{g_s^4}{t^2}\,
 \mathcal S_{\mu\nu}^{(1)}\mathcal S_{(2)}^{\mu\nu}\,\mathcal C,
\end{equation*}
其中
\begin{align*}
 \mathcal S_{\mu\nu}^{(1)}
 &=\frac12\Tr[(\slashed p_3+m_1c)\gamma_\mu(\slashed p_1+m_1c)\gamma_\nu],\\
 \mathcal S_{(2)}^{\mu\nu}
 &=\frac12\Tr[(\slashed p_4+m_2c)\gamma^\mu(\slashed p_2+m_2c)\gamma^\nu].
\end{align*}
颜色部分与质量无关：
\begin{align*}
 \mathcal C
 &=\frac1{N_c^2}\sum_{i,j,k,l}\sum_{a,b}
 (T^a)_{ki}(T^a)_{lj}(T^b)_{ik}(T^b)_{jl}\\
 &=\frac1{N_c^2}\sum_{a,b}\Tr(T^aT^b)\Tr(T^aT^b)\\
 &=\frac{T_F^2(N_c^2-1)}{N_c^2}
 =\frac{N_c^2-1}{4N_c^2}.
\end{align*}
Dirac 迹恒等式给
\begin{align*}
 \mathcal S_{\mu\nu}^{(1)}
 &=2\left[p_{3\mu}p_{1\nu}+p_{3\nu}p_{1\mu}
 -\eta_{\mu\nu}(p_1\!\cdot p_3-m_1^2c^2)\right],\\
 \mathcal S_{(2)}^{\mu\nu}
 &=2\left[p_4^{\mu}p_2^{\nu}+p_4^{\nu}p_2^{\mu}
 -\eta^{\mu\nu}(p_2\!\cdot p_4-m_2^2c^2)\right].
\end{align*}
先在四矢量标量积层面完成张量收缩。把两个方括号分别记为 $A_{\mu\nu}$ 与 $B^{\mu\nu}$，则逐项展开 $A_{\mu\nu}B^{\mu\nu}$ 后，所有
$(p_1\!\cdot p_3)(p_2\!\cdot p_4)$ 的纯迹项相消，得到
\begin{align*}
\mathcal S_{\mu\nu}^{(1)}\mathcal S_{(2)}^{\mu\nu}
=8\Big[&
(p_3\!\cdot p_4)(p_1\!\cdot p_2)
+(p_3\!\cdot p_2)(p_1\!\cdot p_4)\\
&-m_2^2c^2(p_1\!\cdot p_3)
-m_1^2c^2(p_2\!\cdot p_4)
+2m_1^2m_2^2c^4\Big].
\end{align*}
现在才把这些内积换成 Mandelstam 变量。除
\begin{align*}
 2p_1\!\cdot p_2&=s-\Sigma_{12},
&2p_1\!\cdot p_4&=\Sigma_{12}-u,\\
 2p_1\!\cdot p_3&=2m_1^2c^2-t,
&2p_2\!\cdot p_4&=2m_2^2c^2-t,
\end{align*}
外，由四动量守恒还可写成
\begin{equation*}
2p_3\!\cdot p_4=s-\Sigma_{12},
\qquad
2p_3\!\cdot p_2=\Sigma_{12}-u.
\end{equation*}
前两项因而给
\begin{equation*}
\frac14\left[(s-\Sigma_{12})^2+(u-\Sigma_{12})^2\right],
\end{equation*}
而三个显式质量项组合为
\begin{align*}
&-m_2^2c^2\left(m_1^2c^2-\frac t2\right)
-m_1^2c^2\left(m_2^2c^2-\frac t2\right)
+2m_1^2m_2^2c^4\\
&\qquad=\frac{t\Sigma_{12}}2.
\end{align*}
乘回最外层因子 $8$，得到
\begin{equation*}
 \mathcal S_{\mu\nu}^{(1)}\mathcal S_{(2)}^{\mu\nu}
 =2\left[(s-\Sigma_{12})^2+(u-\Sigma_{12})^2+2t\Sigma_{12}\right].
\end{equation*}
与颜色因子相乘即得\eqnrefs{eq:qqprime-M2-massive-dp115}。令 $m_1,m_2\to0$ 后，$\Sigma_{12}\to0$ 且 $s+t+u\to0$，从而得到\eqnrefs{eq:qqprime-M2-final-dp10}。
\end{derivation}

\begin{derivation}[夸克对湮灭的振幅平方]
同时保留初态质量 $m_q$ 与末态质量 $m_Q$。颜色与 Dirac 空间彼此独立，因此
\begin{equation*}
 \overline{|\mathcal M|^2}
 =\frac{g_s^4}{s^2}\,\mathcal C\,\mathcal S.
\end{equation*}
颜色平均为
\begin{align*}
\mathcal C
&=\frac1{N_c^2}
\sum_{i,j,k,l}\sum_{a,b}
(T^a)_{ji}(T^a)_{kl}
(T^b)_{ij}(T^b)_{lk}\\
&=\frac1{N_c^2}
\sum_{a,b}\Tr(T^aT^b)\Tr(T^aT^b)\\
&=\frac{T_F^2(N_c^2-1)}{N_c^2}
=\frac{N_c^2-1}{4N_c^2},
\end{align*}
其中基本表示取 $T_F=1/2$。

自旋平均后的两个 Dirac 迹为
\begin{align*}
 L^{\mu\nu}
 &\equiv\Tr[(\slashed p_2-m_qc)\gamma^\mu
 (\slashed p_1+m_qc)\gamma^\nu]\\
 &=4\left[p_2^\mu p_1^\nu+p_2^\nu p_1^\mu
 -\eta^{\mu\nu}(p_1\!\cdot p_2+m_q^2c^2)\right],\\
 H_{\mu\nu}
 &\equiv\Tr[(\slashed p_3+m_Qc)\gamma_\mu
 (\slashed p_4-m_Qc)\gamma_\nu]\\
 &=4\left[p_{3\mu}p_{4\nu}+p_{3\nu}p_{4\mu}
 -\eta_{\mu\nu}(p_3\!\cdot p_4+m_Q^2c^2)\right].
\end{align*}
于是
\begin{equation*}
 \mathcal S=\frac14L^{\mu\nu}H_{\mu\nu}.
\end{equation*}
先不换成 $s,t,u$，直接把两个对称张量逐项收缩。记
$a\equiv p_1\!\cdot p_2$、$b\equiv p_3\!\cdot p_4$，得到
\begin{align*}
\mathcal S
=8\Big[&
(p_1\!\cdot p_3)(p_2\!\cdot p_4)
+(p_1\!\cdot p_4)(p_2\!\cdot p_3)\\
&+m_Q^2c^2 a+m_q^2c^2 b
+2m_q^2m_Q^2c^4\Big].
\end{align*}
质量壳条件给
\begin{equation*}
a=\frac{s-2m_q^2c^2}{2},
\qquad
b=\frac{s-2m_Q^2c^2}{2}.
\end{equation*}
所以后三项不是独立的三种质量修正，而是精确合并成
\begin{align*}
&m_Q^2c^2\frac{s-2m_q^2c^2}{2}
+m_q^2c^2\frac{s-2m_Q^2c^2}{2}
+2m_q^2m_Q^2c^4\\
&\qquad=\frac{s}{2}(m_q^2+m_Q^2)c^2.
\end{align*}
因此
\begin{equation*}
 \mathcal S
 =8\left[
 (p_1\!\cdot p_3)(p_2\!\cdot p_4)
 +(p_1\!\cdot p_4)(p_2\!\cdot p_3)
 +\frac{s}{2}(m_q^2+m_Q^2)c^2
 \right].
\end{equation*}
令
\begin{equation*}
 \Sigma_{qQ}\equiv(m_q^2+m_Q^2)c^2.
\end{equation*}
由有限质量 Mandelstam 恒等式
\begin{align*}
 2p_1\!\cdot p_3&=\Sigma_{qQ}-t,
 &2p_2\!\cdot p_4&=\Sigma_{qQ}-t,\\
 2p_1\!\cdot p_4&=\Sigma_{qQ}-u,
 &2p_2\!\cdot p_3&=\Sigma_{qQ}-u,
\end{align*}
可得
\begin{equation*}
 \mathcal S
 =2\left[(t-\Sigma_{qQ})^2
 +(u-\Sigma_{qQ})^2
 +2s\Sigma_{qQ}\right].
\end{equation*}
乘回颜色因子得到
\begin{equation*}
 {
 \overline{|\mathcal M|^2}
 =g_s^4\frac{N_c^2-1}{2N_c^2}
 \frac{(t-\Sigma_{qQ})^2+(u-\Sigma_{qQ})^2+2s\Sigma_{qQ}}{s^2},}
\end{equation*}
即正文\eqnrefs{eq:qqbar-to-QQbar-M2-dp10}。外腿质量壳同时给
\begin{equation*}
 s+t+u=2\Sigma_{qQ}.
\end{equation*}
最后沿 $m_q,m_Q\to0$ 的连续质量壳极限得到
\begin{equation*}
 \lim_{m_q,m_Q\to0}\overline{|\mathcal M|^2}
 =g_s^4\frac{N_c^2-1}{2N_c^2}\frac{t^2+u^2}{s^2}.
\end{equation*}
无质量 Born 核因此是有限质量树级结果的严格连续极限。
\end{derivation}

\begin{derivation}[双胶子产生振幅的 Ward 核对]
把两条出射胶子动量记为 $k_1,k_2$，颜色指标为 $a,b$。将三张图共同的 $g_s^2$ 因子提出，并定义约化张量
\begin{equation*}
 \widetilde{\mathcal M}_x^{\mu\nu}
 \equiv g_s^{-2}\mathcal M_x^{\mu\nu},
 \qquad x=t,u,s.
\end{equation*}
对两张夸克交换图，费米子线结构严格写成
\begin{align*}
 \widetilde{\mathcal M}_t^{\mu\nu}
 &=\bar v(p_2)\gamma^\nu T^b
 S(p_1-k_1)\gamma^\mu T^a u(p_1),\\
 \widetilde{\mathcal M}_u^{\mu\nu}
 &=\bar v(p_2)\gamma^\mu T^a
 S(p_1-k_2)\gamma^\nu T^b u(p_1),
\end{align*}
其中 $S(q)$ 表示 Dirac 算符 $S^{-1}(q)=\slashed q-mc$ 的 Feynman 逆核。代数运算先在内部费米子分母非零的开运动学区域进行，所得恒等式再按统一的 $+\ii0$ 处方延拓。

把第一条胶子偏振替换为 $k_1^\mu$。由
\begin{equation*}
 \slashed k_1
 =S^{-1}(p_1)-S^{-1}(p_1-k_1)
\end{equation*}
以及在壳方程 $S^{-1}(p_1)u(p_1)=0$，得到
\begin{align*}
 S(p_1-k_1)\slashed k_1u(p_1)
 &=S(p_1-k_1)
 [S^{-1}(p_1)-S^{-1}(p_1-k_1)]u(p_1)\\
 &=-u(p_1).
\end{align*}
因此 $k_{1\mu}\widetilde{\mathcal M}_t^{\mu\nu}$ 留下一个没有内部传播子的费米子流，颜色顺序为 $T^bT^a$。

对另一种排序，利用四动量守恒 $p_1+p_2=k_1+k_2$，把 $k_1$ 写成靠近反费米子端的两个逆传播子之差，并令逆传播子向左作用到 $\bar v(p_2)$。由 $\bar v(p_2)(\slashed p_2+mc)=0$，得到同一个旋量流，颜色顺序为 $T^aT^b$，符号与上一项相反。两张约化夸克图的纵向收缩之和因而含有
\begin{equation*}
 T^aT^b-T^bT^a=[T^a,T^b]=\ii f^{abc}T^c
\end{equation*}
乘同一个费米子流。阿贝尔理论的生成元是普通数，所以这一对易子严格为零；QCD 的非零对易子要求三胶子图参与同一 Ward 恒等式。

$s$ 道图含夸克颜色流
\begin{equation*}
 J_c^\rho=\bar v(p_2)\gamma^\rho T^c u(p_1)
\end{equation*}
以及由 Yang--Mills 三次作用量得到的三胶子张量 $V_{\rho}{}^{\mu\nu}$。按“全部动量流入”的统一约定收缩 $k_{1\mu}V^{\rho\mu\nu}$，再使用 $k_1^2=k_2^2=0$、$k_2\cdot\varepsilon_2=0$ 与 $q=k_1+k_2$，得到
\begin{equation*}
 k_{1\mu}\widetilde{\mathcal M}_s^{\mu\nu}
 =-k_{1\mu}
 \left(\widetilde{\mathcal M}_t^{\mu\nu}
 +\widetilde{\mathcal M}_u^{\mu\nu}\right).
\end{equation*}
恢复共同因子 $g_s^2$ 后，三图满足
\begin{equation*}
 k_{1\mu}
 \left(\mathcal M_t^{\mu\nu}
 +\mathcal M_u^{\mu\nu}
 +\mathcal M_s^{\mu\nu}\right)=0.
\end{equation*}
交换 $1\leftrightarrow2$ 同样得到第二条 Ward 恒等式。三图抵消依次使用 Dirac 逆传播子的望远镜恒等式、颜色生成元的对易关系和由同一 Yang--Mills 作用量固定的三胶子顶角。
\end{derivation}

\begin{derivation}[三胶子顶角的动量与颜色结构]
Feynman 规范 $\xi=1$ 下，把纯规范二次项与规范固定项分部积分后写成
\begin{equation*}
 \mathcal L_2
 =\frac{\hbar c}{2}\,
 \mathscr C_\mu^a\,\eta^{\mu\nu}\partial^2\mathscr C_\nu^a .
\end{equation*}
因此作用在规范场上的二次微分算符是
\begin{equation*}
 (\mathsf K\mathscr C)_\mu^a
 =\delta^{ab}\eta_{\mu\nu}\partial^2\mathscr C^{b\nu}.
\end{equation*}
采用
\begin{equation*}
 \mathscr C_\mu^a(x)
 =\int\frac{\dd^4p}{(2\pi\hbar)^4}
 \widetilde{\mathscr C}_\mu^a(p)
 \ee^{-\ii p\cdot x/\hbar},
\end{equation*}
有 $\partial^2\mapsto-p^2/\hbar^2$，故
\begin{equation*}
 \mathsf K_{\mu\nu}^{ab}(p)
 =-\frac{p^2}{\hbar^2}\delta^{ab}\eta_{\mu\nu}.
\end{equation*}
Feynman 边界条件把 $p^2$ 替换为 $p^2+\ii0^+$。定义逆核 $\mathsf K_F^{-1}$ 使
\begin{equation*}
 \mathsf K_{\mu\rho}^{ac}(p)
 (\mathsf K_F^{-1})^{cb\rho}{}_{\nu}(p)
 =\delta^{ab}\delta_\mu{}^\nu,
\end{equation*}
于是逐指标求逆得到
\begin{equation*}
 (\mathsf K_F^{-1})_{\mu\nu}^{ab}(p)
 =-\hbar^2\frac{\delta^{ab}\eta_{\mu\nu}}{p^2+\ii0^+}.
\end{equation*}
在统一的约化振幅约定中，玻色 Feynman 线等于 $(\ii/\hbar^2)\mathsf K_F^{-1}$，从而
\begin{equation*}
 D_{\mu\nu}^{ab}(p)
 =\frac{\ii}{\hbar^2}(\mathsf K_F^{-1})_{\mu\nu}^{ab}(p)
 =\frac{-\ii\delta^{ab}\eta_{\mu\nu}}{p^2+\ii0^+},
\end{equation*}
即正文\eqnrefs{eq:qcd-gluon-propagator-dp10}。

三胶子规则来自 Yang--Mills 作用量的三次部分
\begin{equation*}
 \mathcal L_{3g}
 =\hbar c g_s f^{abc}
 (\partial_\mu\mathscr C_\nu^a)
 \mathscr C^{b\mu}\mathscr C^{c\nu}.
\end{equation*}
令三条流入场分别为 $(a,\alpha,p)$、$(b,\beta,q)$、$(c,\gamma,r)$，并约定
\begin{equation*}
 p+q+r=0.
\end{equation*}
三次场都是同一个玻色规范场，因此对作用量取三次泛函导数时，导数场可以对应三条外腿中的任意一条。把所有排列写到统一的颜色次序 $f^{abc}$ 后，三个 Lorentz 系数分别为
\begin{align*}
 \mathscr V_1^{\alpha\beta\gamma}
 &=\eta^{\alpha\beta}(p-q)^\gamma,\\
 \mathscr V_2^{\alpha\beta\gamma}
 &=\eta^{\beta\gamma}(q-r)^\alpha,\\
 \mathscr V_3^{\alpha\beta\gamma}
 &=\eta^{\gamma\alpha}(r-p)^\beta.
\end{align*}
这里每一个动量差都来自两种外腿排列的差：交换相邻两个胶子场时，$f^{abc}$ 变号，而 Fourier 导数把相应流入动量带入同一个 Lorentz 缩并。于是三阶作用量系数给出完整约化顶角
\begin{equation*}
 {
 V_{3g}^{abc,\alpha\beta\gamma}(p,q,r)
 =-g_sf^{abc}
 \left[
 \eta^{\alpha\beta}(p-q)^\gamma
 +\eta^{\beta\gamma}(q-r)^\alpha
 +\eta^{\gamma\alpha}(r-p)^\beta
 \right],}
\end{equation*}
即正文\eqnrefs{eq:qcd-three-gluon-vertex-dp10}。交换任意两条完整外腿，例如 $(a,\alpha,p)\leftrightarrow(b,\beta,q)$，$f^{abc}$ 与方括号同时变号，因此完整顶角保持玻色交换对称。
\end{derivation}

\begin{derivation}[夸克与胶子散射的通道求和]
先取轻夸克极限 $m_q\to0$。这一极限下三张图的相对系数仍由规范不变作用量固定。单张图本身具有规范依赖，因此偏振求和先保留两种横向物理偏振；只有完整振幅相加并满足 Ward 恒等式以后，才可以使用相应的协变化简。

先固定质心系，使每个四动量和物理偏振都能逐分量检查。
令 $Q\equiv\sqrt{s}/2$，并取
\begin{align*}
 p_1^\mu&=Q(1,0,0,1),&
 k_1^\mu&=Q(1,0,0,-1),\\
 p_2^\mu&=Q(1,\sin\theta,0,\cos\theta),&
 k_2^\mu&=Q(1,-\sin\theta,0,-\cos\theta).
\end{align*}
于是
\begin{equation*}
 s=4Q^2,
 \qquad
 t=-2Q^2(1-\cos\theta),
 \qquad
 u=-2Q^2(1+\cos\theta),
 \qquad s+t+u=0.
\end{equation*}
记 $c_\theta\equiv\cos\theta$。入射胶子沿 $-z$ 方向，可选两条实的横向偏振
\begin{equation*}
 \varepsilon_{1x}^\mu=(0,1,0,0),
 \qquad
 \varepsilon_{1y}^\mu=(0,0,1,0),
\end{equation*}
而出射胶子沿 $(-\sin\theta,0,-\cos\theta)$，可选
\begin{equation*}
 \varepsilon_{2\parallel}^\mu=(0,-c_\theta,0,\sin\theta),
 \qquad
 \varepsilon_{2y}^\mu=(0,0,1,0).
\end{equation*}
逐项有 $k_i\cdot\varepsilon_{i\lambda}=0$ 与 $\varepsilon_{i\lambda}^2=-1$。因此这四种偏振组合恰好穷尽两条外胶子的物理态，不含时间样或纵向态。

三张图的 Dirac 核分别定义为
\begin{align*}
 \mathscr A_s^{\mu\nu}
 &=\gamma^\nu\frac{\slashed p_1+\slashed k_1}{s}\gamma^\mu,\\
 \mathscr A_u^{\mu\nu}
 &=\gamma^\mu\frac{\slashed p_1-\slashed k_2}{u}\gamma^\nu,\\
 \mathscr A_t^{\mu\nu}
 &=\frac1t\Gamma^{\rho\mu\nu}(p_1-p_2,k_1,-k_2)\gamma_\rho.
\end{align*}
这里第一条三胶子腿的动量必须是 $p_1-p_2=k_2-k_1$，因为前文三胶子顶角约定三条动量全部流入顶角：
\begin{equation*}
 (p_1-p_2)+k_1-k_2=0.
\end{equation*}
这也是正文\eqnrefs{eq:qcd-qg-amplitude-dp94} 中动量方向的直接检查。

接着计算六个独立的 Dirac--偏振收缩。
对 $x,y\in\{s,u,t\}$ 定义
\begin{equation*}
 K_{xy}(\theta)
 \equiv
 \sum_{r_1,r_2}
 \sum_{\lambda_1,\lambda_2}
 \Bigl[
 \bar u_{r_2}(p_2)\mathscr A_x^{\mu\nu}u_{r_1}(p_1)
 \varepsilon_{1\lambda_1,\mu}
 \varepsilon_{2\lambda_2,\nu}
 \Bigr]
 \Bigl[
 \bar u_{r_2}(p_2)\mathscr A_y^{\alpha\beta}u_{r_1}(p_1)
 \varepsilon_{1\lambda_1,\alpha}
 \varepsilon_{2\lambda_2,\beta}
 \Bigr]^* .
\end{equation*}
利用无质量自旋和
\begin{equation*}
 \sum_r u_r(p)\bar u_r(p)=\slashed p
\end{equation*}
后，每个 $K_{xy}$ 都只是 $4\times4$ Dirac 矩阵的有限迹。含三胶子顶角的项还可以先把 Lorentz 三指标降成一个普通 Dirac 矩阵。令 $q=p_1-p_2=k_2-k_1$，则
\begin{align*}
 &\varepsilon_1^\mu\varepsilon_2^\nu
 \Gamma^{\rho}{}_{\mu\nu}(q,k_1,-k_2)\\
 &\qquad={}
 -2(k_1\!\cdot\varepsilon_2)\varepsilon_1^\rho
 +(\varepsilon_1\!\cdot\varepsilon_2)(k_1+k_2)^\rho
 -2(k_2\!\cdot\varepsilon_1)\varepsilon_2^\rho.
\end{align*}
这一步直接来自三胶子顶角的三项，并只用了
$k_1\cdot\varepsilon_1=k_2\cdot\varepsilon_2=0$ 与 $q=k_2-k_1$。对上面选定的四种物理偏振，它进一步化为
\begin{align*}
 \mathscr A_t^{(x,y)}
 &=-\frac{2Q\sin\theta}{t}\slashed\varepsilon_y,
 &
 \mathscr A_t^{(y,\parallel)}
 &=-\frac{2Q\sin\theta}{t}\slashed\varepsilon_y,\\
 \mathscr A_t^{(y,y)}
 &=-\frac{\slashed k_1+\slashed k_2}{t},
 &
 \mathscr A_t^{(x,\parallel)}
 &=\frac{-2Q\sin\theta\slashed\varepsilon_{1x}
 +c_\theta(\slashed k_1+\slashed k_2)
 -2Q\sin\theta\slashed\varepsilon_{2\parallel}}{t}.
\end{align*}
因此 $t$ 道也已被化成普通的有限伽马矩阵链，不再剩下未计算的三胶子张量。作为第一项示范，$s$ 道平方变成
\begin{align*}
 K_{ss}
 ={}&\sum_{\lambda_1,\lambda_2}
 \Tr\!\left[
 \slashed p_2
 \slashed\varepsilon_2
 \frac{\slashed p_1+\slashed k_1}{s}
 \slashed\varepsilon_1
 \slashed p_1
 \slashed\varepsilon_1
 \frac{\slashed p_1+\slashed k_1}{s}
 \slashed\varepsilon_2
 \right].
\end{align*}
把 $p_1+k_1=(2Q,0,0,0)$ 与上面的四条物理偏振逐一代入，并反复使用
\begin{equation*}
 \{\gamma^\mu,\gamma^\nu\}=2\eta^{\mu\nu}I_4,
 \qquad
 \Tr(\slashed a\slashed b\slashed c\slashed d)
 =4\bigl(a\!\cdot b\,c\!\cdot d-a\!\cdot c\,b\!\cdot d+a\!\cdot d\,b\!\cdot c\bigr),
\end{equation*}
四种偏振组合相加给
\begin{equation*}
 K_{ss}=4(1+c_\theta).
\end{equation*}

再显式算一个含三胶子顶角的干涉项 $K_{st}$。定义 Dirac 共轭矩阵
\begin{equation*}
 \overline{\mathscr A}_t
 \equiv\gamma^0\mathscr A_t^\dagger\gamma^0,
\end{equation*}
则固定一组物理偏振 $(\lambda_1,\lambda_2)$ 时的自旋和为
\begin{equation*}
 K_{st}^{(\lambda_1,\lambda_2)}
 =\Tr\!\left[
 \slashed p_2\,
 \mathscr A_s^{(\lambda_1,\lambda_2)}
 \slashed p_1\,
 \overline{\mathscr A}_t^{(\lambda_1,\lambda_2)}
 \right].
\end{equation*}
把上一段已经化简出的四个 $\mathscr A_t$ 与
$\mathscr A_s=\slashed\varepsilon_2(\slashed p_1+\slashed k_1)\slashed\varepsilon_1/s$
逐一代入，四组偏振分别给
\begin{align*}
 K_{st}^{(x,\parallel)}
 &=\frac{4(1+c_\theta)}{1-c_\theta},&
 K_{st}^{(x,y)}
 &=2(1+c_\theta),\\
 K_{st}^{(y,\parallel)}
 &=2(1+c_\theta),&
 K_{st}^{(y,y)}
 &=\frac{4(1+c_\theta)}{1-c_\theta}.
\end{align*}
这里前、后两项中的 $1/(1-c_\theta)$ 来自 $t=-2Q^2(1-c_\theta)$；分子中的 $Q^2$ 则由两条外动量的 Dirac 迹提供，所以能量尺度消掉，只留下无量纲角函数。四项相加便得到
\begin{align*}
 K_{st}
 &=2(1+c_\theta)
 +2(1+c_\theta)
 +\frac{8(1+c_\theta)}{1-c_\theta}\\
 &=\frac{4(1+c_\theta)(3-c_\theta)}{1-c_\theta}.
\end{align*}
夸克传播子图与三胶子交换图的干涉因此归结为普通 Dirac 迹。物理干涉项还要乘颜色 Gram 元素 $G_{st}<0$，所以 $K_{st}$ 的符号本身不能判断最终干涉是增强还是相消。

剩余四个独立收缩使用同一套迹恒等式得到
\begin{align*}
 K_{uu}&=\frac{8(1+c_\theta^2)}{1+c_\theta},&
 K_{tt}&=\frac{8(1+c_\theta)(c_\theta^2-2c_\theta+5)}{(1-c_\theta)^2},\\
 K_{su}&=4(1-c_\theta),&
 K_{st}&=\frac{4(1+c_\theta)(3-c_\theta)}{1-c_\theta},&
 K_{ut}&=\frac{8(c_\theta^2-c_\theta+2)}{1-c_\theta}.
\end{align*}
由于 $K_{yx}=K_{xy}^*$，这六项已经覆盖三张图平方后的全部九个平方与干涉项。单个 $K_{xy}$ 依赖所选物理偏振基；与颜色结构相干组合并对完整物理振幅求和后，最终结果不依赖这一基底选择。

Lorentz--Dirac 收缩确定后，再独立计算颜色 Gram 矩阵。
把三张图的颜色矩阵记为
\begin{equation*}
 \mathcal C_s=T^bT^a,
 \qquad
 \mathcal C_u=T^aT^b,
 \qquad
 \mathcal C_t=\ii f^{abc}T^c=[T^a,T^b].
\end{equation*}
注意 $\mathcal C_t^\dagger=-\mathcal C_t$，因此干涉项的符号不能靠“颜色因子平方”猜测。定义
\begin{equation*}
 G_{xy}\equiv\sum_{a,b}\Tr_F(\mathcal C_x\mathcal C_y^\dagger).
\end{equation*}
由
\begin{equation*}
 \sum_aT^aT^a=C_FI,
 \qquad
 T^aT^bT^a=\left(C_F-\frac{C_A}{2}\right)T^b,
 \qquad
 \Tr(T^aT^b)=T_F\delta^{ab},
\end{equation*}
逐项得到
\begin{align*}
 G_{ss}=G_{uu}&=T_Fd_AC_F,\\
 G_{su}=G_{us}&=T_Fd_A\left(C_F-\frac{C_A}{2}\right),\\
 G_{tt}&=T_Fd_AC_A,\\
 G_{st}=G_{ts}&=-\frac12T_Fd_AC_A,\\
 G_{ut}=G_{tu}&=+\frac12T_Fd_AC_A,
\end{align*}
其中 $d_A=N_c^2-1$、$T_F=1/2$。例如
\begin{align*}
 G_{st}
 &=\sum_{a,b}\Tr\!\left[T^bT^a(T^bT^a-T^aT^b)\right]\\
 &=T_Fd_A\left(C_F-\frac{C_A}{2}\right)-T_Fd_AC_F
 =-\frac12T_Fd_AC_A,
\end{align*}
负的 $G_{st}$ 与正的 $G_{ut}$ 直接显示非阿贝尔交换图与两种夸克排序具有不同的颜色干涉符号。

最后先在振幅层相干求和，再实施初态平均。
未平均的完整平方振幅为
\begin{align*}
 \sum_{\rm spin,pol,color}|\mathcal M|^2
 =g_s^4\Bigl[&
 G_{ss}K_{ss}+G_{uu}K_{uu}+G_{tt}K_{tt}
 +2G_{su}K_{su}+2G_{st}K_{st}+2G_{ut}K_{ut}
 \Bigr].
\end{align*}
把上面六个 $K_{xy}$ 与六个颜色系数代入，并用
\begin{equation*}
 C_F=\frac{N_c^2-1}{2N_c},
 \qquad C_A=N_c,
 \qquad d_A=N_c^2-1,
\end{equation*}
整理成
\begin{equation*}
 \sum_{\rm spin,pol,color}|\mathcal M|^2
 =g_s^4\frac{d_A}{N_c}
 \frac{(c_\theta^2+2c_\theta+5)
 \left[N_c^2(c_\theta^2+2c_\theta+5)-(1-c_\theta)^2\right]}
 {(1-c_\theta)^2(1+c_\theta)}.
\end{equation*}
初态有 $2N_c$ 个夸克自旋--颜色态和 $2d_A$ 个胶子物理偏振--颜色态，所以平均因子是 $1/(4N_cd_A)$。再使用
\begin{align*}
 s^2+u^2&=4Q^4(c_\theta^2+2c_\theta+5),\\
 t^2&=4Q^4(1-c_\theta)^2,\\
 su&=-8Q^4(1+c_\theta),
\end{align*}
就得到
\begin{equation*}
 \overline{|\mathcal M(qg\to qg)|^2}
 =g_s^4(s^2+u^2)
 \left[
 \frac1{t^2}-\frac{C_F}{C_A}\frac1{su}
 \right],
\end{equation*}
即正文\eqnrefs{eq:qcd-qg-M2-dp94}。这里 $1/t^2$ 与 $1/(su)$ 是六类运动学收缩和六类颜色 Gram 元素相干组合后的结果；它们不能分别归属于某一张 Feynman 图的“概率”。
\end{derivation}

\begin{derivation}[胶子散射的颜色与偏振求和]
这一过程若直接用 $-\eta_{\mu\nu}$ 展开四条外腿的偏振求和，会产生很长的 Lorentz 指标收缩，而且容易把“偏振代数”和“颜色代数”混在一起。这里改用物理螺旋度基底：先把每个非零螺旋度振幅算成三个颜色张量的线性组合，再单独完成一般 $SU(N_c)$ 的颜色求和。这样每一步都只处理一种数学结构。

先固定质心系运动学与外胶子的物理偏振。
令 $Q\equiv\sqrt{s}/2$，取度规 $\eta_{\mu\nu}=\operatorname{diag}(1,-1,-1,-1)$，并写
\begin{align*}
 k_1^\mu&=Q(1,0,0,1),&
 k_2^\mu&=Q(1,0,0,-1),\\
 k_3^\mu&=Q(1,\sin\theta,0,\cos\theta),&
 k_4^\mu&=Q(1,-\sin\theta,0,-\cos\theta).
\end{align*}
于是
\begin{equation*}
 t=-\frac{s}{2}(1-\cos\theta),
 \qquad
 u=-\frac{s}{2}(1+\cos\theta),
 \qquad s+t+u=0.
\end{equation*}
四个外胶子都只取横向物理偏振。以右螺旋度为例，可选
\begin{align*}
 \varepsilon_R^\mu(k_1)&=\frac1{\sqrt2}(0,1,\ii,0),&
 \varepsilon_R^\mu(k_2)&=\frac1{\sqrt2}(0,-1,\ii,0),\\
 \varepsilon_R^{*\mu}(k_3)&=\frac1{\sqrt2}(0,-\cos\theta,\ii,\sin\theta),&
 \varepsilon_R^{*\mu}(k_4)&=\frac1{\sqrt2}(0,\cos\theta,\ii,-\sin\theta).
\end{align*}
左螺旋度由上式作复共轭得到。逐一检查可得
\begin{equation*}
 k_i\cdot\varepsilon_{R,L}(k_i)=0,
 \qquad
 \varepsilon_{R,L}^*(k_i)\cdot\varepsilon_{R,L}(k_i)=-1,
\end{equation*}
所以这里从一开始就没有引入纵向或时间样偏振。

为缩短颜色指标，定义
\begin{equation*}
 \mathcal C_s\equiv f^{abe}f^{cde},
 \qquad
 \mathcal C_t\equiv f^{ace}f^{bde},
 \qquad
 \mathcal C_u\equiv f^{ade}f^{bce}.
\end{equation*}
Jacobi 恒等式给出
\begin{equation*}
 \mathcal C_s-\mathcal C_t+\mathcal C_u=0.
\end{equation*}

随后逐图计算独立的螺旋度振幅。
先算 $g_Rg_R\to g_Rg_R$。为避免与整体系数的无关相位纠缠，下表列出费曼规则给出的 $\ii\mathcal M$。三行交换图分别含两个三胶子顶角和一个 $1/s$、$1/t$、$1/u$ 传播子，最后一行为四胶子接触顶角。把上面的四个偏振矢量逐分量代入三胶子顶角
\begin{equation*}
 \Gamma^{\mu\nu\rho}(p,q,r)
 =\eta^{\mu\nu}(p-q)^\rho
 +\eta^{\nu\rho}(q-r)^\mu
 +\eta^{\rho\mu}(r-p)^\nu,
 \qquad p+q+r=0,
\end{equation*}
即可得到
\begin{center}
\begin{tabular}{c|c}
图 & $\ii\mathcal M(g_Rg_R\to g_Rg_R)/(\ii g_s^2)$\\ \hline
$s$ 道 & $-\mathcal C_s\cos\theta$\\[2pt]
$t$ 道 & $\displaystyle
 \mathcal C_t\frac{19+7\cos\theta-11\cos^2\theta+\cos^3\theta}
 {4(1-\cos\theta)}$\\[8pt]
$u$ 道 & $\displaystyle
 \mathcal C_u\frac{19-7\cos\theta-11\cos^2\theta-\cos^3\theta}
 {4(1+\cos\theta)}$\\[8pt]
四胶子接触项 & $\displaystyle
 -\mathcal C_s\cos\theta
 -\frac{\mathcal C_t}{4}(3+2\cos\theta-\cos^2\theta)
 -\frac{\mathcal C_u}{4}(3-2\cos\theta-\cos^2\theta)$
\end{tabular}
\end{center}
这里每一行只是有限个四维分量的乘加。以 $s$ 道为例，传播子分母为 $(k_1+k_2)^2=s$；两个三胶子顶角分别与
$\varepsilon_R(k_1)\varepsilon_R(k_2)$ 和
$\varepsilon_R^*(k_3)\varepsilon_R^*(k_4)$ 收缩。利用
$k_i\cdot\varepsilon_i=0$ 后，分子中的 $Q^2=s/4$ 恰好抵消传播子的 $1/s$，留下表中的 $-\mathcal C_s\cos\theta$。$t$、$u$ 道保留的 $1/(1\mp\cos\theta)$ 正是后来前向、后向交换极点的来源；接触项没有传播子，因此只给角度多项式。

把四行相加，并用 $\mathcal C_s=\mathcal C_t-\mathcal C_u$，分子多项式发生大规模相消：
\begin{align*}
 \frac{\ii\mathcal M_{RR\to RR}}{\ii g_s^2}
 &=4\left(
 \frac{\mathcal C_t}{1-\cos\theta}
 +\frac{\mathcal C_u}{1+\cos\theta}
 \right)\\
 &=-2\left(
 \mathcal C_t\frac{s}{t}
 +\mathcal C_u\frac{s}{u}
 \right).
\end{align*}
这一步显示了四胶子接触项的作用：它本身没有交换极点，却与三个交换图的多项式部分相消，使完整振幅只剩满足 Ward 恒等式与交叉结构的组合。

其余非零螺旋度振幅由同一分量计算或交叉变换得到。取三个独立通道可写成
\begin{align*}
 \mathcal M_{RR\to RR}
 &=-2g_s^2s\left(\frac{\mathcal C_t}{t}+\frac{\mathcal C_u}{u}\right),\\
 \mathcal M_{LR\to LR}
 &=2g_s^2u\left(\frac{\mathcal C_t}{t}+\frac{\mathcal C_s}{s}\right),\\
 \mathcal M_{LR\to RL}
 &=-2g_s^2t\left(\frac{\mathcal C_s}{s}-\frac{\mathcal C_u}{u}\right).
\end{align*}
整体相位不影响平方振幅。宇称变换给出另外三个等模振幅；其余螺旋度组合在树级为零。因此螺旋度求和只需把这三个独立结果平方后乘以 $2$；因子 $2$ 来自宇称相关的另外三个等模通道。

螺旋度核确定后，计算一般 $SU(N_c)$ 的颜色 Gram 矩阵。
记伴随表示维数
\begin{equation*}
 d_A=N_c^2-1.
\end{equation*}
首先，三个颜色张量的平方范数相同。例如
\begin{align*}
 \sum_{a,b,c,d}\mathcal C_s^2
 &=\sum_{e,e'}
 \left(\sum_{a,b}f^{abe}f^{ab e'}\right)
 \left(\sum_{c,d}f^{cde}f^{cd e'}\right)\\
 &=C_A^2\sum_{e,e'}\delta^{ee'}\delta^{ee'}
 =C_A^2d_A.
\end{align*}
令伴随表示生成元为 $(F^a)_{bc}=-\ii f^{abc}$。由
\begin{equation*}
 \Tr(F^aF^b)=C_A\delta^{ab},
 \qquad
 \sum_aF^aF^bF^a=\frac{C_A}{2}F^b,
\end{equation*}
可得一个代表性的交叉颜色积
\begin{align*}
 \sum_{a,b,c,d}\mathcal C_s\mathcal C_t
 &=\sum_{e,e'}\Tr(F^eF^{e'}F^eF^{e'})\\
 &=\frac{C_A}{2}\sum_{e'}\Tr(F^{e'}F^{e'})
 =\frac12C_A^2d_A.
\end{align*}
再结合 Jacobi 关系 $\mathcal C_s-\mathcal C_t+\mathcal C_u=0$ 及外颜色指标的置换对称性，得到完整 Gram 矩阵
\begin{equation*}
 G_{xy}\equiv\sum_{a,b,c,d}\mathcal C_x\mathcal C_y
 =C_A^2d_A
 \begin{pmatrix}
 1&\tfrac12&-\tfrac12\\
 \tfrac12&1&\tfrac12\\
 -\tfrac12&\tfrac12&1
 \end{pmatrix}_{xy},
 \qquad x,y\in\{s,t,u\}.
\end{equation*}
负的 $G_{su}$ 并不表示负概率；它只记录两个颜色张量在颜色向量空间中的夹角，最终概率由完整振幅的 Gram 二次型给出。

最后完成颜色求和、螺旋度求和与初态平均。
利用上面的 Gram 矩阵，第一个独立螺旋度通道给出
\begin{align*}
 \sum_{\rm color}|\mathcal M_{RR\to RR}|^2
 &=4g_s^4C_A^2d_A\,s^2
 \left(
 \frac1{t^2}+\frac1{u^2}+\frac1{tu}
 \right).
\end{align*}
同理
\begin{align*}
 \sum_{\rm color}|\mathcal M_{LR\to LR}|^2
 &=4g_s^4C_A^2d_A\,u^2
 \left(
 \frac1{t^2}+\frac1{s^2}+\frac1{ts}
 \right),\\
 \sum_{\rm color}|\mathcal M_{LR\to RL}|^2
 &=4g_s^4C_A^2d_A\,t^2
 \left(
 \frac1{s^2}+\frac1{u^2}+\frac1{su}
 \right).
\end{align*}
宇称伙伴再给相同贡献，因此全螺旋度、全颜色求和为
\begin{equation*}
 \sum_{\rm hel,color}|\mathcal M|^2
 =8g_s^4C_A^2d_A\,\mathcal S(s,t,u),
\end{equation*}
其中
\begin{align*}
 \mathcal S(s,t,u)={}&
 s^2\left(\frac1{t^2}+\frac1{u^2}+\frac1{tu}\right)
 +t^2\left(\frac1{s^2}+\frac1{u^2}+\frac1{su}\right)\\
 &+u^2\left(\frac1{s^2}+\frac1{t^2}+\frac1{st}\right).
\end{align*}
利用 $s+t+u=0$，例如
\begin{equation*}
 -\frac{2su}{t^2}
 =\frac{s^2+u^2}{t^2}-1,
\end{equation*}
以及循环置换，可把它化成
\begin{equation*}
 \mathcal S(s,t,u)
 =2\left(
 3-\frac{su}{t^2}-\frac{st}{u^2}-\frac{tu}{s^2}
 \right).
\end{equation*}
每个初态胶子有 $d_A$ 个颜色和两个物理螺旋度，所以初态简并度为 $4d_A^2$。最终
\begin{align*}
 \overline{|\mathcal M|^2}
 &=\frac{1}{4d_A^2}
 \sum_{\rm hel,color}|\mathcal M|^2\\
 &=\frac{4g_s^4C_A^2}{d_A}
 \left(
 3-\frac{su}{t^2}-\frac{st}{u^2}-\frac{tu}{s^2}
 \right),
\end{align*}
再代入 $d_A=N_c^2-1$，即得到正文\eqnrefs{eq:qcd-gg-M2-dp94}。
\end{derivation}

\begin{derivation}[相同夸克散射的交换干涉]
把两张图写为
\begin{align*}
\mathcal M_t&=\frac{g_s^2}{t}
[\bar u_3\gamma^\mu T^a u_1]
[\bar u_4\gamma_\mu T^a u_2],\\
\mathcal M_u&=\frac{g_s^2}{u}
[\bar u_4\gamma^\mu T^a u_1]
[\bar u_3\gamma_\mu T^a u_2].
\end{align*}
相同末态费米子交换给 $\mathcal M=\mathcal M_t-\mathcal M_u$，所以
\begin{equation*}
\overline{|\mathcal M|^2}
=\overline{|\mathcal M_t|^2}
+\overline{|\mathcal M_u|^2}
-2\Re\overline{\mathcal M_t\mathcal M_u^*}.
\end{equation*}
前两项直接由不同味结果得到
\begin{equation*}
\overline{|\mathcal M_t|^2}
=g_s^4\frac{N_c^2-1}{2N_c^2}\frac{s^2+u^2}{t^2},
\qquad
\overline{|\mathcal M_u|^2}
=g_s^4\frac{N_c^2-1}{2N_c^2}\frac{s^2+t^2}{u^2}.
\end{equation*}
真正的新信息在交叉项。颜色平均可写成一个单迹：
\begin{align*}
\mathcal C_{tu}
&=\frac1{N_c^2}\sum_{a,b}\Tr(T^aT^bT^aT^b).
\end{align*}
利用
\begin{equation*}
T^aT^bT^a=\left(C_F-\frac{C_A}{2}\right)T^b
=-\frac{1}{2N_c}T^b,
\end{equation*}
以及 $\sum_b\Tr(T^bT^b)=T_F(N_c^2-1)$，得到
\begin{equation*}
\mathcal C_{tu}
=-\frac{N_c^2-1}{4N_c^3}.
\end{equation*}
自旋交叉项不能拆成两个独立迹，而是形成一个八 $\gamma$ 链。先按完备关系把四条外腿的自旋求和，注意复共轭振幅会反转第二条费米子链的矩阵次序，于是
\begin{equation*}
\mathcal S_{tu}
=\frac14\Tr\!\left(
\slashed p_3\gamma^\mu\slashed p_1\gamma^\nu
\slashed p_4\gamma_\mu\slashed p_2\gamma_\nu
\right).
\end{equation*}
这里不能把迹“按视觉上相邻的四个矩阵”拆开。先把中间五个 $\gamma$ 矩阵作为整体收缩。四维 Clifford 代数给出
\begin{equation*}
\gamma^\mu\gamma^\alpha\gamma^\nu\gamma^\beta\gamma_\mu
=-2\gamma^\beta\gamma^\nu\gamma^\alpha,
\end{equation*}
因此
\begin{align*}
\mathcal S_{tu}
&=-\frac12\Tr\!\left(
\slashed p_3\slashed p_4\gamma^\nu
\slashed p_1\slashed p_2\gamma_\nu
\right).
\end{align*}
第二次收缩只剩四个普通四矢量。利用
\begin{equation*}
\gamma^\nu\slashed a\slashed b\gamma_\nu=4a\!\cdot b,
\qquad
\Tr(\slashed a\slashed b)=4a\!\cdot b,
\end{equation*}
逐步得到
\begin{align*}
\mathcal S_{tu}
&=-\frac12\,4(p_1\!\cdot p_2)\Tr(\slashed p_3\slashed p_4)\\
&=-8(p_1\!\cdot p_2)(p_3\!\cdot p_4).
\end{align*}
对无质量 $2\to2$ 运动学，$2p_1\!\cdot p_2=s$ 且 $2p_3\!\cdot p_4=s$，故
\begin{equation*}
{\mathcal S_{tu}=-2s^2.}
\end{equation*}
这一步说明干涉项只留下 $s^2$ 并不是由 $s+t+u=0$ 猜出的，而是八 $\gamma$ 链在 Clifford 代数下真正收缩后的结果。
于是
\begin{equation*}
-2\Re\overline{\mathcal M_t\mathcal M_u^*}
=-2\frac{g_s^4}{tu}\,\mathcal C_{tu}\mathcal S_{tu}
=-g_s^4\frac{N_c^2-1}{N_c^3}\frac{s^2}{tu},
\end{equation*}
与两个平方项相加即得到\eqnrefs{eq:qcd-identical-qq-m2-dp95}。注意颜色迹本身为负，而 Fermi 交换又给振幅相对负号；最终干涉符号必须在这两层结构全部保留后才能确定。
\end{derivation}
